\subsection*{Android to Node-RED Migration Guide}

\textbf{Prompt:}
\begin{lstlisting}[language={}, breaklines=true]
I need a technical guide on migrating an IoT companion app's frontend from Android (Kotlin/Jetpack Compose) to a Node-RED UI Builder web interface (HTML/CSS/Vanilla JS).

Using the provided context (parts of our Android codebase and the current project README), please outline:

1. A mapping of Android concepts to Web/Node-RED equivalents.
2. Architectural guidelines for structuring the HTML, CSS, and JS files modularly.
3. Simple code templates demonstrating how to parse incoming telemetry JSON cleanly and dispatch structured commands back via `uibuilder.send()` over WebSockets.

Focus on the general patterns rather than our game's business logic, and explain the key differences in state management between Jetpack Compose and modern Vanilla JavaScript.
\end{lstlisting}

\textbf{Context:}
\begin{lstlisting}[language={}, breaklines=true]
Relevant parts of the Android Source Code, including basic styling, constants, modelview, JSON commands, MQTT handling, and the README describing the current architecture and design patterns.
\end{lstlisting}

\textbf{Answer:}
\begin{lstlisting}[language={Java}, breaklines=true]
Migrating an IoT companion interface from a native Android application to a Node-RED UI Builder web dashboard involves moving from a statically typed, declarative JVM environment to a dynamic, event-driven web runtime.

1. CONCEPT MAPPING
+-------------------------+---------------------------------+
| Android (JVM / Kotlin)  | Web / UI Builder Equivalent     |
+-------------------------+---------------------------------+
| SharedPreferences       | LocalStorage API                |
| Jetpack Compose State   | Stateful JS Object + DOM Binder |
| MqttAsyncClient         | uibuilder WebSocket connection  |
| Coroutine Delay/Timeout | setTimeout() / clearTimeout()   |
+-------------------------+---------------------------------+

2. ARCHITECTURAL STRUCTURING

UI Builder uses a modular directory structure under 'src/'. To mirror Android's separation of concerns, split the codebase as follows:
- index.html: Declarative markup structure (the layouts).
- index.css: Visual properties and CSS Custom Properties (the theme).
- index.js: Entry point coordinating events and state changes.
- app/types.js: Holds factories and default configurations for fallback handling.
- app/command.js: Assembles and dispatches commands over WebSockets.
- app/render.js: Dynamic rendering engine updating the DOM tree.

3. BASIC GENERIC PATTERNS

A. Saving and Loading Connection Profiles (LocalStorage):

// Saving
const credentials = { host: "localhost", port: 1883 };
localStorage.setItem('my_app_creds', JSON.stringify(credentials));

// Loading
const stored = localStorage.getItem('my_app_creds');
const myCreds = stored ? JSON.parse(stored) : { host: "default", port: 1883 };

B. Parsing Telemetry securely with Fallbacks:

function parseTelemetry(rawPayload) {
    let json = {};
    try {
        json = (typeof rawPayload === 'string') ? JSON.parse(rawPayload) : rawPayload;
    } catch (e) {
        console.error("Payload error", e);
    }
    return {
        status: json.status || "unknown",
        value: typeof json.value === 'number' ? json.value : 0.0
    };
}

C. Dispatching Outbound Commands via UI Builder:

function sendCommand(uibuilder, action, parameters) {
    const payload = {
        command: action,
        timestamp: new Date().toISOString(),
        payload: parameters
    };
    uibuilder.send({
        payload: payload,
        topic: "my_app/command"
    });
}

4. STATE MANAGEMENT DIFFERENCES

- Jetpack Compose uses a compiler plugin to automatically track reads/writes to mutable states (e.g. mutableStateOf) and trigger micro-recompositions of corresponding UI trees.
- In Vanilla JS, state updates do not automatically re-render elements. To manage state, store parameters inside a plain JavaScript object and call explicit render/update functions whenever a mutating event occurs.
\end{lstlisting}